%MB9T5
\begin{vidu}%CTST
Trong thí nghiệm sóng dừng trên sợi dây AB dài 120 cm với đầu B tự do, đầu A được kích thích để thực hiện dao động với biên độ nhỏ. Ngoài đầu A, trên dây xuất hiện thêm một nút. 
\begin{itemize}
\item [1.] Tần số sóng là 12,5 Hz. Tốc độ truyền sóng  là ......... 
\item [2.] Độ lệch giữa hai tần số gần nhau nhất để có sóng dừng trên dây  là ......... 
\item [3.] Tần số bé nhất để có sóng dừng trên dây là ......... 
\item [4.] Để có thêm một nút sóng thì tần số sóng khi này phải bằng  ......... 
\item [5.] Giữ cố định đầu B của dây. Tần số bé nhất để có sóng dừng lúc này là ......... 
\item [6.] Khi giữ đầu B cố định. Cho tần số thay đổi từ 0 tới 40 Hz. Số lần tạo sóng dừng trên dây  là ......... 
\end{itemize}
\loigiai{
\begin{itemize}
\item[1.] Tốc độ truyền sóng: $\ell=\dfrac{3\lambda}{4}=\dfrac{3v}{4f}\Rightarrow v=\dfrac{4\ell f}{3}=20\ m/s$. 
\item [2.] Độ lệch giữa hai tần số liên tiếp là $f_{k+1}-f_k=\left(2(k+1)+1\right)\dfrac{v}{4\ell}-\left(2k+1\right)\dfrac{v}{4\ell}=\dfrac{v}{2\ell}=\dfrac{25}{3}\ Hz$
\item [3.] Tần số bé nhất để có sóng dừng trên dây là: $f_{\min}=\dfrac{v}{4\ell}=\dfrac{25}{6}\ Hz$
\item[4.] Để có thêm một nút sóng thì $\ell=\dfrac{5\lambda}{4}=\dfrac{5v}{4f}\Rightarrow f=\dfrac{5v}{4\ell}=\dfrac{125}{6}\ Hz$.
\item [5.] Khi giữ cố định đầu B. Tần số bé nhất tạo sóng dừng là: $f_{\min}=\dfrac{v}{2\ell}=\dfrac{25}{3}\ Hz$
\item [6.] Ta có $0< f=k\dfrac{v}{2\ell}=k\dfrac{25}{3}\le 40\Rightarrow 0< k\le 4,8 \Rightarrow $ có 4 lần tạo sóng dừng
\item 
\item 
\item 
\end{itemize}
}
\end{vidu} %\dotlineEX{19}
\loai{Biết hai tần số liên tiếp, tính tần số nhỏ nhất tạo ra sóng dừng}
\begin{vidu}%[3P2-9.1K][Loai1]
Sóng dừng trên dây dài một đầu cố định, một đầu tự do. Hai tần số gần nhau nhất cùng tạo ra sóng dừng trên dây là 175 Hz và 225 Hz. Tần số nhỏ nhất tạo ra sóng dừng trên dây đó là
\choice
{50 Hz}
{\True 25 Hz}
{75 Hz}
{100 Hz}
\loigiai{
\begin{itemize}
\item Hiệu hai tần số liên tiếp để dây có sóng dừng là: $\Delta f= 225 – 175 = 50$ Hz.
\item   Tần số nhỏ nhất để tạo ra sóng dừng một đầu cố định, một đầu tự do là: $f_{\min }=\dfrac{\Delta f}{2}=25$ Hz.
\end{itemize}
}
\end{vidu} %\dotlineEX{3}
\loai{Chưa biết hai đầu cố định hay tự do}
\begin{vidu}%[3P2-12.1K][Loai1]}
Một dây đàn hồi tạo sóng dừng với ba tần số liên tiếp là 75 Hz; 125 Hz và 175 Hz. Biết dây thuộc loại hai đầu cố định hoặc có một đầu cố định, đầu kia tự do và vận tốc truyền sóng trên đây là 400 m/s. Chiều dài của dây là 
\choice
{8 m}
{5 m}
{\True 4 m}
{7 m}
\loigiai{
\begin{itemize}
\item Nếu trường hợp dây thuộc loại hai đầu cố định thì các tần số tiếp theo sẽ là $f_0$, 2$f_0$, 3$f_0$... và $\Delta f = f_0$
\item Nếu dây một đầu hở, một đầu cố định các tần số liên tiếp là $f_0$, 3$f_0$, 5 $f_0$, 7$f_0$..và $\Delta f = 2f_0$
\item Từ 3 tần số liên tiếp mà đề bài cho $\Rightarrow \Delta f = 50$ Hz.
\item Xét trường hợp $\Delta f = f_0 = 50$ Hz. Dễ thấy $\dfrac{75}{50} = 1,5$ là số không nguyên nên dây đàn này không thể có hai đầu cố định được.
\item Xét trường hợp $\Delta f = 2f_0 = 50 \Rightarrow f_0 = 25$ Hz. Dễ thấy $\dfrac{75}{25} = 3$ là số nguyên nên dây đàn này một đầu tự do, một đầu cố định.
\item $f_0 = \dfrac{v}{4\ell} \Rightarrow\ell = \dfrac{v}{4f_0} = \dfrac{400}{4.25} = 4\ m$.
\end{itemize}
}
\end{vidu} %\dotlineEX{4}
\loai{Hai tần số liên tiếp cùng tạo ra sóng dừng. Tìm tốc độ truyền sóng}
\begin{vidu}%[3P2-9.1K][Loai1]
Một sợi dây căng giữa hai điểm cố định cách nhau 75 cm. Người ta tạo sóng dừng trên dây. Hai tần số gần nhau nhất cùng tạo ra sóng dừng trên dây là 150 Hz và 200 Hz. Tốc độ truyền sóng trên dây là
\choice
{7,5 m/s}
{300 m/s}
{225 m/s}
{\True 75 m/s}
\loigiai{
\begin{itemize}
\item Hiệu hai tần số liên tiếp để dây có sóng dừng là: $f_0=\dfrac{v}{2\ell }= 200 – 150 = 50\ Hz \Rightarrow v = 75$ m/s.
\end{itemize}
}
\end{vidu} %\dotlineEX{4}

\loai{Điều kiện kẹp tần số- Tìm tần số}
\begin{vidu}%[3P2-9.1K][Loai1]
Người ta tạo sóng dừng trên một dây AB dài 0,8 m. Đầu B cố định, đầu A gắn với cần rung với tần số f  ($35 \ Hz \leq f$). Biết tốc độ truyền sóng trên dây bằng 8 m/s. Tần số sóng \textit{có thể} bằng
\choice
{38 Hz}
{\True 40 Hz}
{42 Hz}
{36 Hz}
\loigiai{
\begin{itemize}
\item Ta có $AB = k\dfrac{v}{2f}$. Với $ f \geq 35$.
\item Ta thấy $k=8$ thì $f=40$ Hz nên B đúng.
\end{itemize}
}
\end{vidu} %\dotlineEX{2}
\loai{Điều kiện kẹp tần số- Tìm số lần tạo sóng dừng}
\begin{vidu}%[3P2-9.1K][Loai1]
Một sợi dây đàn hồi dài 1,2 m được treo lơ lửng lên một cần rung. Cần có thể rung theo phương ngang với tần số thay đổi được từ 100 Hz đến 125 Hz. Tốc độ truyền sóng trên dây là 6 m/s. Biết rằng khi có sóng dừng, đầu nối với cần rung là nút sóng. Trong quá trình thay đổi tần số rung của cần, có thể tạo ra được bao nhiêu lần sóng dừng trên dây? 
\choice
{\True 10 lần}
{4 lần}
{5 lần}
{12 lần}
\loigiai{
\begin{itemize}
\item Chiều dài sợi dây: $\ell = \left(k + \dfrac{1}{2}\right) \dfrac{\lambda}{2} = \left(k + \dfrac{1}{2}\right)\dfrac{v}{2f}\Rightarrow f = \left(k + \dfrac{1}{2}\right)\dfrac{v}{2\ell} = 2,5(k + 0,5)$.
\item Mà $100\leq f \leq 125 \Rightarrow 39,5\leq k \leq 49,5 \Rightarrow$ Có 10 giá trị k thỏa mãn.
\end{itemize}
}
\end{vidu} %\dotlineEX{5}
\loai{Điều kiện kẹp tần số - Tìm số nút, bụng}
\begin{vidu}%[3P2-9.1K][Loai1]
\nguon{MH - 2018} Một sợi dây dài 2 m với hai đầu cố định, đang có sóng dừng. Sóng truyền trên dây với tốc độ 20 m/s. Biết rằng tần số của sóng truyền trên dây có giá trị trong khoảng từ 11 Hz đến 19 Hz. Tính cả hai đầu dây, số nút sóng trên dây là
\choice
{5}
{3}
{\True 4}
{2}
\loigiai{
\begin{itemize}
\item Ta có điều kiện sóng dừng trên hai đầu dây cố định: $\ell =k\dfrac{\lambda}{2}=k\dfrac{v}{2f}\Rightarrow f=k\dfrac{v}{2\ell}$.
\item Mặt khác:
$11\ Hz<f<19\ Hz\Rightarrow 11<k\dfrac{v}{2\ell}<19\Rightarrow 2,2<k<3,8\Rightarrow k=3$.
\item Số nút sóng $=k+1=4$ nút.
\end{itemize}
}
\end{vidu} %\dotlineEX{5}
\loai{Số nút, bụng thay đổi}
\begin{vidu}%[3P2-9.1K][Loai1]
Một sợi dây đàn hồi, hai đầu cố định có sóng dừng. Khi tần số sóng trên dây là 30 Hz thì trên dây có 3 bụng sóng. Muốn trên dây có 4 bụng sóng thì phải
\choice
{\True tăng tần số thêm 10 Hz}
{tăng tần số thêm 30 Hz}
{giảm tần số đi 10 Hz}
{giảm tần số còn $\dfrac{20}{3}$ Hz}
\loigiai{
\begin{itemize}
\item Ta có: $\ell = \dfrac{k\lambda}{2} = \dfrac{kv}{2f}\Rightarrow f = \dfrac{kv}{2\ell}$.
\item Lập tỉ số: $\dfrac{f_2}{f_1} = \dfrac{k_2}{k_1} = \dfrac{4}{3}\Rightarrow f_2 = \dfrac{4f_1}{3} = 40\ Hz $.
\item Vậy phải tăng tần số thêm 10 Hz.
\end{itemize}
}
\end{vidu} %\dotlineEX{4}

\loai{Số nút, bụng thay đổi kèm tần số}
\begin{vidu}%[3P2-9.1K][Loai1]
Một dây AB hai đầu cố định. Khi dây rung với tần số f thì trên dây có 4 bó sóng. Khi tần số tăng thêm 10 Hz thì trên dây có 5 bó sóng, vận tốc truyền sóng trên dây là 10 m/s. Chiều dài và tần số rung của dây là
\choice
{\True $\ell=50$ cm, $f=40$ Hz}
{$\ell=40$ cm, $f=50$ Hz}
{$\ell=5$ cm, $f=50$ Hz}
{$\ell=50$ cm, $f=50$ Hz}
\loigiai{
\begin{itemize}
\item Sóng dừng với hai đầu cố định ta có: $\ell = \dfrac{k \lambda}{2} = \dfrac{kv}{2f} \Rightarrow f = \dfrac{kv}{2\ell}.$
\item Ban đầu: $f = \dfrac{4v}{2\ell}\quad (1)$
\item Lúc sau: $f + 10 =  \dfrac{5v}{2\ell}\quad (2)$
\item Từ $(1)$ và $(2)$ ta có $\ell=50$ cm, $f=40$ Hz.
\end{itemize}
}
\end{vidu} %\dotlineEX{5}
\loai{Thay đổi tần số lượng nhỏ nhất, không thay đổi tính chất 2 đầu}
\begin{vidu}%[3P2-9.1K][Loai1]
Trên một sợi dây đàn hồi dài 50 cm buông thẳng đứng, đầu trên cố định, đầu dưới tự do đang có một sóng dừng với tần số xác định. Để trên dây tiếp tục hình thành sóng dừng thì phải thay đổi tần số kích thích một lượng nhỏ nhất bằng 20 Hz. Tốc độ truyền sóng trên dây bằng
\choice
{17,5 m/s}
{37,5 m/s}
{27,5 m/s}
{\True 20 m/s}
\loigiai{
\begin{itemize}
\item Ta có: $\ell = \left(n + \dfrac{1}{2}\right)\dfrac{v}{2\ell} = (2n + 1)\dfrac{v}{4f}\Rightarrow f = \left(2n + 1\right)\dfrac{v}{4\ell} = \left(2n + 1\right)f_0$ với $f_0 = \dfrac{v}{4\ell}$.
\item $ \Delta f = 2f_0 = 20 \Rightarrow f_0 = 10\ Hz\Rightarrow v = 4\ell.f_0 = 4.0,5.10 = 20\ m/s$. 
\end{itemize}
}
\end{vidu} %\dotlineEX{5}
\loai{Thay đổi đầu cố định tự do. Tần số thay đổi lượng nhỏ nhất}
\begin{vidu}%[3P2-9.1G][Loai1]
Một sợi dây đàn hồi dài 90 cm có một đầu gắn với nguồn dao động, một đầu tự do. Khi dây rung với tần số 10 Hz thì trên dây xuất hiện sóng dừng với 5 bụng trên dây. Nếu đầu tự do của đầu dây được giữ cố định và tốc độ truyền sóng trên dây không đổi thì phải thay đổi tần số rung của dây một lượng nhỏ nhất là bao nhiêu để tiếp tục có sóng dừng trên dây
\choice
{\True $\dfrac{10}{9}$ Hz}
{$\dfrac{10}{11}$ Hz}
{$\dfrac{11}{9}$ Hz}
{12 Hz}
\loigiai{
\begin{itemize}
\item Lúc đầu có sóng dừng một đầu cố định một đầu tự do:
$10\ Hz=\left( 2. 5-1 \right)\dfrac{v}{4\ell }\Rightarrow \dfrac{v}{4\ell }=\dfrac{10}{9}$ Hz $\Rightarrow \dfrac{v}{2\ell }=\dfrac{20}{9}$ Hz.
\item Cố định đầu tự do, nếu vẫn giữ nguyên tần số thì: $10=n\dfrac{v}{2\ell }\Rightarrow n=4,5$.
\item Nếu tăng tần số lượng nhỏ nhỏ nhất sẽ có sóng dừng 5 bụng $\Rightarrow f_5 = 5\dfrac{v}{2\ell }=\dfrac{100}{9}\ Hz$. 
\item Tăng tần số thêm lượng là: $\dfrac{100}{9}-10=\dfrac{10}{9}\ Hz$.
\item Nếu tần tần số lượng nhỏ nhỏ nhất sẽ có sóng dừng 4 bụng $\Rightarrow f_4$ = $4\dfrac{v}{2\ell }=\dfrac{80}{9}\ Hz$.
\item Giảm tần số đi lượng là: $10-\dfrac{80}{8}=\dfrac{10}{9}\ Hz$.
\item Vậy cả hai trường hợp tần số đã thay đổi lượng $\dfrac{10}{9}\ Hz$.
\end{itemize}
}
\end{vidu} %\dotlineEX{4}
\loai{Thay đổi đầu cố định tự do. Tìm tần số, tốc độ truyền sóng}
\begin{vidu}%[3P2-9.1K][Loai1]
Trên thanh AB đàn hồi với hai đầu cố định đang có một sóng dừng với tần số bằng 50 Hz. Buông đầu B cho tự do thì để trên thanh vẫn có sóng dừng cần phải rút ngắn độ dài thanh đi một đoạn nhỏ nhất bằng 10 cm. Tốc độ truyền sóng trên thanh bằng
\choice
{5 m/s}
{10 m/s}
{15 m/s}
{\True 20 m/s}
\loigiai{
\begin{itemize}
\item Khi hai đầu dây cố định thì $\ell = n\dfrac{\lambda}{2}$.
\item Sau khi để đầu B tự do, để thanh vẫn có sóng dừng thì cần rút ngắn đi một lượng nhỏ nhất $\Delta \ell_{\min} = \dfrac{\lambda}{4} = 10 \Rightarrow \lambda = 40\ cm $.
\item Tốc độ truyền sóng: $v = \lambda.f = 40.50 = 2000\ cm /s = 20\ m/s$.
\end{itemize}
}
\end{vidu} %\dotlineEX{5}
\loai{Biện luận nghiệm nguyên}
\begin{vidu}%[3P2-9.1T][Loai1]
Một dây đàn hồi căng ngang, một đầu cố định, một đầu tự do. Thấy hai tần số tạo ra sóng dừng trên dây là 2964 Hz và 4940 Hz. Biết tần số nhỏ nhất tạo ra sóng dừng nằm trong khoảng từ 216 Hz đến 524 Hz. Với tần số nằm trong khoảng từ 8 kHz đến 11 kHz, có bao nhiêu tần số tạo ra sóng dừng?
\choice
{6}
{7}
{8}
{\True 5}
\loigiai{
\begin{itemize}
\item Luôn có: $f=\left( 2n+1 \right)\dfrac{v}{4\ell }=\left( 2n+1 \right). {{f}_{\min }}$; Điều kiện: $216\ Hz < f_{\min} < 524\ Hz\quad (*)$
\item Tần số 2964 Hz, giả sử có $m +1$ bụng sóng $\Rightarrow 2964 = (2m + 1)\Rightarrow f_{\min} = \dfrac{2964}{2m+1}\quad (1)$. 
\item Từ $(*) \Rightarrow 5,7 < 2m + 1 < 13,7$.
\item Tần số 4940 Hz, giả sử có $k+1$ bụng sóng $\Rightarrow 4940 = (2k + 1)\Rightarrow f_{\min} =\dfrac{4940}{2k-1}\quad (2)$ 
\item Từ $(*) \Rightarrow  9,4 < 2k +1 < 22,9$.
\item Chia từng vế $(1)$ với $(2) \Rightarrow \dfrac{2m+1}{2k+1}=\dfrac{3}{5}=\dfrac{9}{15}=\dfrac{15}{25}... $ 
\item So sánh với điều kiện của m và k $\Rightarrow \dfrac{2m-1}{2k-1}= \dfrac{9}{15}$ là cặp nghiệm phù hợp. \\ 
$\Rightarrow$$f_{\min}$ = $\dfrac{988}{3}\ Hz$ $\Rightarrow$  Tần số để xảy ra sóng dừng phải là: $f = \dfrac{988}{3}\left( 2n-1 \right)$.
\item   Mặt khác theo đề: $8000 < \dfrac{988}{3}\left( 2n+1 \right) < 11000 \Rightarrow 11,6 < n < 16,2$.
\item   $n = 12,\ 13,\ 14,\ 15,\ 16:$ có 5 tần số trong khoảng đó có thể tạo ra sóng dừng.

\end{itemize}
}
\end{vidu} %\dotlineEX{15}

